%% ============================================================================
%% Fichier  : annexes/Ann_B_outils.tex
%% Titre    : Référentiel des outils
%% Version  : 1.0 -- novembre 2026
%% ============================================================================

\chapter{Référentiel des Outils}
\label{ann:outils}

Les outils cités sont des exemples. Le cours privilégie les outils simples
(tableur, langage de programmation) qui obligent à comprendre le calcul.

\begin{tableaucours}{L{3.8cm}L{5.8cm}L{2.6cm}}{%
    \tableauHeader{Outil} &
    \tableauHeader{Usage} &
    \tableauHeader{Chapitre, lab}
}
Tableur & Registre, matrice, cotation, tableau de bord & Ch.~\ref{chap:matrices}, ARGC-LAB-01 \\
Cartographie de l'écosystème & Menace des parties prenantes (EBIOS RM) & Ch.~\ref{chap:ebios}, ARGC-LAB-02 \\
Estimation à trois points & Facteurs FAIR, moyenne PERT & Ch.~\ref{chap:quantitatif}, ARGC-LAB-03 \\
Python ou tableur & Simulation de Monte-Carlo, centiles & Ch.~\ref{chap:simulation}, ARGC-LAB-04 \\
Tableau de bord & Indicateurs clés de risque et seuils & Ch.~\ref{chap:traiter}, ARGC-LAB-05 \\
Plan d'échantillonnage & Taille d'échantillon, limite supérieure & Ch.~\ref{chap:conduire}, ARGC-LAB-06 \\
Registre des incidents & Classement, délais, escalade & Ch.~\ref{chap:incidents}, ARGC-LAB-07 \\
Grille d'analyse d'impact & MTPD, RTO, RPO, dépendances & Ch.~\ref{chap:bia}, ARGC-LAB-08 \\
Main courante & Chronologie, décisions, actions & Ch.~\ref{chap:cellule}, ARGC-LAB-09 \\
Modèles de messages & Message d'attente, questions et réponses & Ch.~\ref{chap:communiquer}, ARGC-LAB-10 \\
\end{tableaucours}
